\section{Discussion}

\subsection{Three claims}

\textbf{Nonnegative source weighting is insufficient for the convention
class studied here.} Theorem~\ref{thm:nonneg} proves it; the ablation in
\S6.3 shows it. Anything that multiplies messages by nonnegative scalars
inherits the failure, no matter how the scalars are estimated.

\textbf{Signed scalars are necessary but not sufficient once conventions
depend on context.} Theorem~\ref{thm:scalar} makes this precise: the
effective sign of a source's contribution is a function of context, and no
context-independent scalar can represent a function. The signed-to-context-conditioned
gap in the ablation is the empirical fingerprint.

\textbf{Explicit context conditioning recovers much of the lost
performance, but Bayesian inference over the latent generative convention
provides an additional advantage from uncertainty marginalisation.}
\textsc{ContextConditioned} reaches $0.730$ against Bayes's $0.758$, and
the EM ablation shows the gap decomposes into $2.6$ points from
point-estimate versus posterior-averaging and $4.9$ points from hard
decoding. Theorem~\ref{thm:consistency} shows the posterior concentrates
on the observational-equivalence class of the truth as $K \to \infty$;
Experiment~C traces that trajectory empirically. The diagnostic
\textsc{MetaAttn+Dec}, given privileged decoded features, reaches $0.822$
against Bayes's $0.758$, so the learned models' bottleneck is
representational. Two further stress tests sharpen the claim. First, a
discriminative source${}\times{}$context meta-learner
(\textsc{CtxLogReg}, \S6.3) with the same support information reaches
only $0.637$, below both \textsc{ContextConditioned} and Bayes:
expressiveness without the right inductive bias does not help at $K=10$.
Second, Experiment~H breaks the model's independence assumption with
correlated source errors and finds Bayes hurt no more than the baselines
or the oracle. With matched support information and modest capacity,
building in the generative structure beats learning it from scratch.

\subsection{What the results do not say}

The benchmark is synthetic: a controlled binary task, involutive
conventions, uniform priors, deterministic contexts. Real systems have
continuous messages, non-involutive mappings, and structured contexts.
None of our theorems depend on binarity in an essential way, but the
experiments do not test beyond it. Correlated sources are now tested
(Experiment~H finds no disproportionate harm), but the correlation form
is one convenient choice, not an exhaustive study.

The convention family is finite, and the Bayesian method assumes the
candidate set is known. Open-world systems would need to grow the set
adaptively. Experiment~G probes the incomplete-family case and finds smooth
rather than catastrophic degradation, but a principled open-world extension
is future work, not this paper.

Posterior concentration is moderate at small $K$ and sharpens with support
size, just as Theorem~\ref{thm:consistency} predicts.

\subsection{Design implications}

Two practical lessons. First, \emph{do not aggregate raw messages across
heterogeneous sources}: averaging messages from agents with different
conventions can do worse than chance. Second, \emph{calibrate before
aggregating}: keep a support set per source and combine calibrated
log-odds rather than raw messages. Reliability heterogeneity is about
\emph{how often} a source is right; convention heterogeneity is about
\emph{how} it maps readings to messages. A source can be perfectly
reliable and completely unreadable.
